\subsubsection{Recovering Probabilities from Expectation Values}

In the previous sections, we have seen that the expectation value of an observable is the weighted average of its possible measurement outcomes. But, is possible to go in the opposite direction? That is, \hl{if we know the expectation value, can we determine the probabilities of the individual outcomes?} Let's explore how to got \definition{Probabilities from the Expectation Value}.

\highspace
Let's consider the Pauli-Z observable as an example.
\begin{equation*}
    Z
    =
    \begin{bmatrix}
        1 & 0 \\
        0 & -1
    \end{bmatrix}
\end{equation*}
With (proof at \autopageref{sec:measuring-ket-with-the-pauli-z-observable}):
\begin{itemize}
    \item Eigenvalue $+1$ associated with the state $\ket{0}$;
    \item Eigenvalue $-1$ associated with the state $\ket{1}$.
\end{itemize}
Let:
\begin{itemize}
    \item $p_0$ be the probability of measuring the state $\ket{0}$;
    \item $p_1$ be the probability of measuring the state $\ket{1}$.
\end{itemize}
Since Pauli-Z is a valid observable, the probabilities must satisfy:
\begin{equation*}
    p_0 + p_1 = 1
\end{equation*}
Therefore, the expectation value of $Z$ is:
\begin{equation*}
    \langle Z \rangle
    =
    (+1)p_0
    +
    (-1)p_1
    =
    p_0 - p_1
\end{equation*}
Where $+1$ and $-1$ are the eigenvalues of $Z$ corresponding to the states $\ket{0}$ and $\ket{1}$, respectively. Instead, $p_0$ and $p_1$ are the probabilities of obtaining those states upon measurement.

\highspace
\textcolor{Green3}{\faIcon[regular]{lightbulb}} Now, we are interested in the opposite direction: given the expectation value $m = \langle Z \rangle$ (we change the notation just to be general), can we determine $p_0$ and $p_1$? The answer is yes, and we can derive the formulas for $p_0$ and $p_1$ in terms of $m$ (or $\langle Z \rangle$). We have two relationships:
\begin{align*}
    p_0 + p_1 &= 1 \\
    p_0 - p_1 &= m
\end{align*}
We can \hl{solve this system of equations to express $p_0$ and $p_1$ in terms of $m$, and thus recover the probabilities from the expectation value.} Let's do that step by step.

\newpage

\noindent
Let's try to solve this system of equations to express $p_0$ and $p_1$ in terms of $m$:
\begin{align*}
    \begin{cases}
        p_0 + p_1 = 1 \\
        p_0 - p_1 = m
    \end{cases}
    &=
    \begin{cases}
        p_0 = 1 - p_1 \\
        p_1 = - m + p_0
    \end{cases} \\
    &=
    \begin{cases}
        p_0 = 1 - p_1 \\
        p_1 = - m + (1 - p_1)
    \end{cases} \\
    &=
    \begin{cases}
        p_0 = 1 - p_1 \\
        2p_1 = - m + 1
    \end{cases} \\
    &=
    \begin{cases}
        p_0 = 1 - p_1 \\
        p_1 = \frac{1 - m}{2}
    \end{cases} \\
    &=
    \begin{cases}
        p_0 = 1 - \frac{1 - m}{2} \\
        p_1 = \frac{1 - m}{2}
    \end{cases} \\
    &=
    \begin{cases}
        p_0 = \frac{2 - (1 - m)}{2} \\
        p_1 = \frac{1 - m}{2}
    \end{cases} \\
    &=
    \begin{cases}
        p_0 = \frac{1 + m}{2} \\
        p_1 = \frac{1 - m}{2}
    \end{cases}
\end{align*}
So, we have:
\begin{equation}\label{eq:probabilities-from-expectation}
    p_0
    =
    \frac{1 + m}{2}
    \quad\text{and}\quad
    p_1
    =
    \frac{1 - m}{2}
\end{equation}
\textcolor{Green3}{\faIcon{question-circle} \textbf{Why is this important?}} This result shows that the \hl{expectation value of the Pauli-Z observable contains enough information to reconstruct the complete probability distribution of the measurement outcomes.} In other words, by measuring only the average value $m = \langle Z \rangle$, we can determine:
\begin{itemize}
    \item The probability $p_0$ of obtaining the state $\ket{0}$;
    \item The probability $p_1$ of obtaining the state $\ket{1}$.
\end{itemize}

\highspace
\begin{flushleft}
    \textcolor{Red2}{\faIcon{exclamation-triangle} \textbf{Not always true in general}}
\end{flushleft}
For a general observable with more than two outcomes, one expectation value is usually \textbf{not enough}. In general, \hl{to determine $n$ unknown probabilities, we need $n$ independent equations.} Normalization provides one equation, so we typically need additional information such as: more expectation values, higher moments (e.g., $\left\langle O^2 \right\rangle$), or other constraints to fully determine the probability distribution.

\highspace
However, if an observable has only two distinct outcomes $a$ and $b$, then the same idea works. If:
\begin{itemize}
    \item Outcome $a$ occurs with probability $p$,
    \item Outcome $b$ occurs with probability $1-p$,
\end{itemize}
Then the expectation value is:
\begin{equation}
    \left\langle O \right\rangle
    =
    a \cdot p + b \cdot (1-p)
    \xRightarrow{\text{solve for } p}
    p
    =
    \frac{\left\langle O \right\rangle - b}{a - b}
\end{equation}
It finds \textbf{the probability of the first outcome $a$}, and the \textbf{probability of the second outcome $b$} can be found by normalization:
\begin{equation}
    1 - p
    =
    1 - \frac{\left\langle O \right\rangle - b}{a - b}
    =
    \frac{a - \left\langle O \right\rangle}{a - b}
\end{equation}
For example, for the Pauli-Z observable, we have $a = +1$ and $b = -1$, so we can recover the probabilities as:
\begin{equation*}
    p_0 = P(+1) = \frac{\langle Z \rangle - (-1)}{(+1) - (-1)} = \frac{\langle Z \rangle + 1}{2}
\end{equation*}
\begin{equation*}
    p_1 = P(-1) = \frac{(+1) - \langle Z \rangle}{(+1) - (-1)} = \frac{1 - \langle Z \rangle}{2}
\end{equation*}